% VI38-DETAILED-PROBLEM-06
\begin{problem}[VI.38.P6: Quadratic Veronese embedding]\label{prob:vi38-06}
Prove from graded rings that \(\nu_2:\mathbb P^1\to\mathbb P^2\), \([s:t]\mapsto[s^2:st:t^2]\), is a closed immersion and compute its ideal.
\end{problem}

\ifdefined\FullProblemDossiers
\begin{solution}
Let
\[
S=k[s,t],
\qquad
S^{(2)}=\bigoplus_{r\ge0}S_{2r}.
\]
The regraded ring \(S^{(2)}\) is generated in degree one by
\(s^2,st,t^2\). Hence
\[
\theta:k[Y_0,Y_1,Y_2]\twoheadrightarrow S^{(2)},
\qquad
Y_0\mapsto s^2,
\quad Y_1\mapsto st,
\quad Y_2\mapsto t^2
\]
is a surjective graded map.  The quadratic relation
\[
Q=Y_0Y_2-Y_1^2
\]
lies in \(\ker\theta\).

We prove that it generates the kernel.  Using
\(Y_0Y_2\equiv Y_1^2\pmod{(Q)}\), every monomial can be reduced to a monomial
containing not both \(Y_0\) and \(Y_2\); thus it has one of the forms
\[
Y_0^aY_1^b
\quad\text{or}\quad
Y_1^bY_2^c.
\]
Their images are respectively
\[
s^{2a+b}t^b,
\qquad
s^bt^{b+2c}.
\]
Within a fixed total degree these exponent pairs determine the reduced monomial
uniquely (the two families meet only in powers of \(Y_1\)). Therefore distinct
reduced monomials have distinct images in the monomial basis of \(k[s,t]\). No
nonzero linear combination of them lies in the kernel, so
\[
\boxed{\ker\theta=(Y_0Y_2-Y_1^2)}.
\]

The quotient map induces a closed immersion
\[
\Proj S^{(2)}\hookrightarrow\mathbb P^2
\]
with image
\[
V_+(Y_0Y_2-Y_1^2).
\]
Finally, Veronese invariance gives
\[
\Proj S\cong\Proj S^{(2)},
\]
and under this isomorphism the degree-one generators \(s^2,st,t^2\) give precisely
\[
\boxed{\nu_2([s:t])=[s^2:st:t^2]}.
\]
Thus \(\nu_2\) is a closed immersion onto the projective conic with ideal
\((Y_0Y_2-Y_1^2)\).
\end{solution}
\fi
